% ========================================================
% BÀI TẬP TOÁN 5 - KẾT NỐI TRI THỨC VỚI CUỘC SỐNG
% Bài 34: Ôn tập đo lường
% Định dạng: BTTeX Part III - Trắc nghiệm điền số ngắn
% ========================================================

% Câu 1
\begin{ex}
  Một ki-lô-mét bằng bao nhiêu mét?
  \shortans{1000}
  \loigiai{
    $1\text{ km} = 1\,000\text{ m}$.
  }
\end{ex}

% Câu 2
\begin{ex}
  Một mét bằng bao nhiêu xăng-ti-mét?
  \shortans{100}
  \loigiai{
    $1\text{ m} = 100\text{ cm}$.
  }
\end{ex}

% Câu 3
\begin{ex}
  Một tấn bằng bao nhiêu ki-lô-gam?
  \shortans{1000}
  \loigiai{
    $1\text{ tấn} = 1\,000\text{ kg}$.
  }
\end{ex}

% Câu 4
\begin{ex}
  Một tạ bằng bao nhiêu yến?
  \shortans{10}
  \loigiai{
    $1\text{ tạ} = 10\text{ yến}$.
  }
\end{ex}

% Câu 5
\begin{ex}
  Một héc-ta bằng bao nhiêu mét vuông?
  \shortans{10000}
  \loigiai{
    $1\text{ ha} = 10\,000\text{ m}^2$.
  }
\end{ex}

% Câu 6
\begin{ex}
  Một ki-lô-mét vuông bằng bao nhiêu héc-ta?
  \shortans{100}
  \loigiai{
    $1\text{ km}^2 = 100\text{ ha}$.
  }
\end{ex}

% Câu 7
\begin{ex}
  Đổi $3\text{ km } 25\text{ m}$ thành bao nhiêu mét?
  \shortans{3025}
  \loigiai{
    $3\,000 + 25 = 3\,025\text{ m}$.
  }
\end{ex}

% Câu 8
\begin{ex}
  Đổi $4\text{ tấn } 5\text{ tạ}$ thành bao nhiêu tạ?
  \shortans{45}
  \loigiai{
    $40 + 5 = 45\text{ tạ}$.
  }
\end{ex}

% Câu 9
\begin{ex}
  Đổi $2\text{ tạ } 50\text{ kg}$ thành bao nhiêu ki-lô-gam?
  \shortans{250}
  \loigiai{
    $200 + 50 = 250\text{ kg}$.
  }
\end{ex}

% Câu 10
\begin{ex}
  Đổi $5\text{ m}^2\text{ } 45\text{ dm}^2$ thành bao nhiêu đề-xi-mét vuông?
  \shortans{545}
  \loigiai{
    $500 + 45 = 545\text{ dm}^2$.
  }
\end{ex}

% Câu 11
\begin{ex}
  Đổi $8\text{ dm}^2\text{ } 6\text{ cm}^2$ thành bao nhiêu xăng-ti-mét vuông?
  \shortans{806}
  \loigiai{
    $800 + 6 = 806\text{ cm}^2$.
  }
\end{ex}

% Câu 12
\begin{ex}
  Đổi $2\text{ giờ } 15\text{ phút}$ thành bao nhiêu phút?
  \shortans{135}
  \loigiai{
    $120 + 15 = 135\text{ phút}$.
  }
\end{ex}

% Câu 13
\begin{ex}
  Đổi $3\text{ phút } 40\text{ giây}$ thành bao nhiêu giây?
  \shortans{220}
  \loigiai{
    $180 + 40 = 220\text{ giây}$.
  }
\end{ex}

% Câu 14
\begin{ex}
  Một ngày đêm có bao nhiêu giờ?
  \shortans{24}
  \loigiai{
    Một ngày đêm có 24 giờ.
  }
\end{ex}

% Câu 15
\begin{ex}
  Một năm không nhuận có bao nhiêu ngày?
  \shortans{365}
  \loigiai{
    Năm không nhuận có 365 ngày.
  }
\end{ex}

% Câu 16
\begin{ex}
  Một năm nhuận có bao nhiêu ngày?
  \shortans{366}
  \loigiai{
    Năm nhuận có 366 ngày.
  }
\end{ex}

% Câu 17
\begin{ex}
  Thế kỉ thứ XXI bắt đầu từ năm $2001$. Thế kỉ này có bao nhiêu năm?
  \shortans{100}
  \loigiai{
    Mỗi thế kỉ có 100 năm.
  }
\end{ex}

% Câu 18
\begin{ex}
  Một con voi cân nặng khoảng $3$ tấn. Con voi đó nặng bao nhiêu tạ?
  \shortans{30}
  \loigiai{
    $3\text{ tấn} = 30\text{ tạ}$.
  }
\end{ex}

% Câu 19
\begin{ex}
  Một bao xi măng nặng $50\text{ kg}$. Hỏi $20$ bao xi măng như thế nặng bao nhiêu tấn?
  \shortans{1}
  \loigiai{
    $50 \times 20 = 1\,000\text{ kg} = 1\text{ tấn}$.
  }
\end{ex}

% Câu 20
\begin{ex}
  Đổi $4{,}5\text{ m}$ thành bao nhiêu đề-xi-mét?
  \shortans{45}
  \loigiai{
    $4{,}5 \times 10 = 45\text{ dm}$.
  }
\end{ex}

% Câu 21
\begin{ex}
  Đổi $2{,}75\text{ kg}$ thành bao nhiêu gam?
  \shortans{2750}
  \loigiai{
    $2{,}75 \times 1\,000 = 2\,750\text{ g}$.
  }
\end{ex}

% Câu 22
\begin{ex}
  Đổi $0{,}5\text{ ha}$ thành bao nhiêu mét vuông?
  \shortans{5000}
  \loigiai{
    $0{,}5 \times 10\,000 = 5\,000\text{ m}^2$.
  }
\end{ex}

% Câu 23
\begin{ex}
  Một người đi bộ với vận tốc $4\text{ km/h}$. Trong $15$ phút người đó đi được bao nhiêu mét?
  \shortans{1000}
  \loigiai{
    $15\text{ phút} = \frac{1}{4}\text{ giờ}$. Quãng đường: $4 \times \frac{1}{4} = 1\text{ km} = 1\,000\text{ m}$.
  }
\end{ex}

% Câu 24
\begin{ex}
  Một bình chứa $1{,}5$ lít nước ngọt. Rót đều vào các cốc $250\text{ ml}$. Rót được bao nhiêu cốc?
  \shortans{6}
  \loigiai{
    $1\,500 : 250 = 6$ cốc.
  }
\end{ex}

% Câu 25
\begin{ex}
  Một bể chứa nước có thể tích $3\text{ m}^3$. Bể đó chứa được tối đa bao nhiêu nghìn lít nước (biết $1\text{ dm}^3 = 1\text{ l}$)?
  \shortans{3}
  \loigiai{
    $3\text{ m}^3 = 3\,000$ lít = 3 nghìn lít.
  }
\end{ex}
